\documentclass[10pt,a4paper]{amsart}
\usepackage[margin=19mm]{geometry}
\usepackage{amsmath,amssymb,mathtools}
\usepackage[scr=rsfs,cal=euler]{mathalfa}
\usepackage{xfrac}
\usepackage[unicode,hidelinks,bookmarks=false]{hyperref}
\newcommand{\ie}{i.e.\ }
\newcommand{\cf}{cf.\ }
\newcommand{\resp}{respectively\ }
\newcommand{\Subsection}{\subsection}
\providecommand{\nobreakdash}{\nobreak\hbox{-}}
\setlength{\emergencystretch}{2em}
\multlinegap0pt
\usepackage{bm}
\usepackage{bbm}
\usepackage{dsfont}
\usepackage{xspace}
\usepackage{cancel}
\usepackage{mathtools}
\usepackage{amsthm}
\makeatletter
\@ifundefined{theorem}{\newtheorem{theorem}{Theorem}}{}
\@ifundefined{Lemma}{\newtheorem{Lemma}[theorem]{Lemma}}{}
\makeatother
\title[Universal Deformations of a Curve and a Differential]{Universal Deformations of a Curve and a Differential}
\author{Emma Carberry, Martin Ulrich Schmidt}
\date{Source: arXiv:2211.11442v2 (CC BY 4.0)}
\begin{document}
\maketitle

\noindent\emph{Source and licence.} Universal Deformations of a Curve and a Differential. Version arXiv:2211.11442v2 (CC BY 4.0), distributed under the Creative Commons Attribution 4.0 International licence, as recorded in the arXiv licence metadata for this exact version and shown on the abstract page https://arxiv.org/abs/2211.11442v2. Full rights notice: licenses/arxiv-2211.11442-CC-BY-4.0.txt.\par\smallskip

\noindent\textit{Editorial scope.} Counted unit: the Lemma whose source label is \texttt{l1} in the article (source0.tex lines 703--708, source label \texttt{l1}), delivered together with the article's own complete original proof of it (source0.tex lines 709--728, one \texttt{proof} environment of 20 lines). No mathematics is written, repaired or re-derived: statement and proof are the article's own text line for line. Required context retained uncounted: eq:mor (lines 569--571). Every definition, display and label of those lines is retained unchanged. Omitted from this package and not redistributed: 1--568, 572--702, 729--1310. Nothing inside a retained excerpt is omitted or altered: every retained body is delivered exactly as the article's lines are written, so no omission receipt of any kind accompanies this package. Citations: the retained excerpts carry no citation command at all, so no bibliography entry of the article is transcribed and no reference is redistributed. Reference adaptation: none was needed and none was made; every reference inside the delivered unit resolves inside it, so no unresolved reference or citation is produced. Printed slips kept literal: no printed word, symbol, hypothesis or number of the counted statement or of its proof was altered. Layout and class: the article's own class is \texttt{amsart} with packages including verbatim, amscd, epsfig, amsmath, amsthm, amstext, amssymb, hyperref, extpfeil, tikz-cd; the wrapper is the lane's pinned \texttt{amsart} editorial wrapper at 10pt on A4, which restates the theorem-like environments the retained text needs and adds the article's own macro definitions that the retained text needs (none), with extra wrapper packages (bm, bbm, dsfont, xspace, cancel, mathtools). The delivered numbering is the unit's own. Assets: no retained span contains a figure environment, a graphics-inclusion command, a PDF-page inclusion, a table or a TikZ picture; no image file is copied into this package, the package declares no asset dependency and no image file is redistributed with it. Licence: version arXiv:2211.11442v2 of this article is distributed under the Creative Commons Attribution 4.0 International licence, as recorded in the arXiv licence metadata for this exact version and shown on the abstract page https://arxiv.org/abs/2211.11442v2. A full rights notice is delivered at licenses/arxiv-2211.11442-CC-BY-4.0.txt. Characters: every retained excerpt is pure ASCII, so the delivered engine, XeLaTeX with its default Latin Modern text font, reports no missing character.
\begin{align}\label{eq:mor}
 H(x,y,s)\Tilde{F}(x,u(x,y,s),\phi(s))&=F(x,y,s).
\end{align}

\begin{Lemma}\label{l1}
The  $x$-isomorphism classes of the infinitesimal deformations of a space germ $ V (f) $ with $f\in\mathbb{C}\{x,y\},\, f(0,0)=0$ are in one-to-one correspondence with 
\begin{align}\label{eq:quotient}
\mathbb{C}\{x,y\}\big/\langle f,\tfrac{\partial f}{\partial y}\rangle.
\end{align}
\end{Lemma}

\begin{proof}
 An infinitesimal deformation of $ V (f) $ is given by $ V (F) $ for $F(\epsilon)=f+\epsilon g$ and $g\in\mathbb{C}\{x,y\}$. % We shall relax our terminology and refer directly to $ F (\epsilon) $
Suppose we are given two infinitesimal  deformations of $ V (f) $, defined by $F_1(\epsilon)=f+\epsilon g_1$ and $F_2(\epsilon)=f+\epsilon g_2$ respectively. An $ x $-morphism  between them has the form
\[
\begin{array}{cclcl}
V (f) &\hookrightarrow& V (F_1) &\twoheadrightarrow& \mathcal I_\epsilon\\
\|&&\hspace{1mm}\downarrow\psi&&\downarrow\phi\\
V (f) &\hookrightarrow& V (F_2) &\twoheadrightarrow& \mathcal I_\epsilon.
\end{array}
\]
Note however that any automorphism $\phi$ of $\mathcal I_\epsilon $ can be lifted to an automorphism of the infinitesimal deformation. Hence for the purpose of classifying $x$-isomorphism classes of the infinitesimal deformations, we may assume that $\phi$ is the identity. An $x$-morphism for which $\phi $ is the identity is explicitly described, as in ~\eqref{eq:mor}, by  $u,H\in\mathbb{C}\{x,y\}$ such that
$$(1+\epsilon H(x,y))(f(x,y+\epsilon u(x,y))+
\epsilon g_2(x,y+\epsilon u(x,y))=f(x,y)+\epsilon g_1(x,y)$$
holds in $\mathbb{C}\{x,y\}\otimes\Pi$. The $\epsilon^0$-term and the $\epsilon^1$-term yield
\begin{align*}
f(x,y)&=f(x,y),&
H(x,y)f(x,y)+u(x,y)\frac{\partial f}{\partial y}(x,y)+g_2(x,y)&=g_1(x,y).
\end{align*}
We conclude that the two infinitesimal deformations $F_1(\epsilon)=f+\epsilon g_1$ and $F_2(\epsilon)=f+\epsilon g_2$ are $x$-isomorphic if and only if $g_1=g_2$ in $\mathbb{C}\{x,y\}/\langle f,\frac{\partial f}{\partial y}\rangle$.
\end{proof}

\end{document}
